\section{Setup}\label{sec:setup}

\subsection{Data and the reasoning probe}
Models read raw bytes (vocabulary 256, context 512), so cross-entropy converts directly to bits per byte
(bpb). Text comes from a 256\,MiB slice of a ClimbMix-derived web corpus~\citep{diao2025climb}; a
separate shard gives 8\,MiB of held-out text, of which we score 384 windows ($\approx$197k bytes).

The reasoning probe is \emph{variable tracking}~\citep{hsieh2024ruler} made dense. A line declares
single-letter variables that hold a digit or point at an earlier variable, in random order:
\begin{center}\texttt{q=4;k=q;x=7;b=k;y=x;?b=}\,$\rightarrow$\,\texttt{4}\end{center}
Exactly one chain has $n$ pointer hops; distractor variables are digits or point at earlier
distractors, never deeper than $n$. Answering requires $n$ serial lookups, so deeper chains need more
sequential computation, which is what loops are supposed to provide. In training, 8 of every 32 rows are
such lines with $n\in\{0,\dots,6\}$ and every variable is queried, so each answer byte is a target. We
evaluate exact match of the argmax digit for the deepest variable on 400 fixed problems per depth:
\textbf{in-distribution} (ID) at $n\in\{2,4,6\}$ and \textbf{out-of-distribution} (OOD) at
$n\in\{8,10\}$, deeper than anything seen in training. Chance is 10\%. Unit tests check that every
generated answer and depth is correct.

\subsection{Model}
The baseline is a looped transformer: one prelude block, two shared core blocks applied $R=4$ times,
and one coda block (effective depth $1+2R+1=10$), with RMSNorm, RoPE and SwiGLU, $d=384$, six heads and
7.28M parameters. The prelude output $e$ is re-injected at every loop: $h_{r}=\mathrm{Core}(h_{r-1}+e)$.
Experiment switches, all in one editable file:
\begin{itemize}
\item \textbf{Depth attention} (\texttt{DEPTH\_ATTN=$K$}, after ARLoop~\citep{yao2026arloop}): before each
  loop, every token attends over its own last $K$ loop states,
  $x \leftarrow x + W_o\,\mathrm{softmax}(q(x)^\top k(S)/\sqrt{d_k})\,v(S)$, with $W_o$ initialised
  to zero.
\item \textbf{Adaptive halting} (\texttt{HALTING}, after ACT and Mixture-of-Recursions
  \citep{graves2016adaptive,bae2025mixture}): a router gives each token a halting probability per loop.
  Training mixes loop states by their halting distribution and adds a ponder cost $\lambda\,\mathbb{E}[R]$;
  evaluation freezes a token once its probability passes a threshold, and we report the realised mean
  loops per token.
\item \textbf{KV reuse} (\texttt{KV\_SHARE}, after \citep{huang2026looped,licakova2026kvreplay}): loops
  after the first attend with the first loop's keys and values, so the core contributes one layer's KV
  per block instead of $R$; the cache shrinks from 15 to 6\,KB per token.
\item \textbf{History injection} (\texttt{HISTORY}, \texttt{LOOP\_EMBED}, after
  \citep{zhuang2026recurrence}): add a (gated) mean of earlier loop states and a sinusoidal loop-index
  embedding, both through zero-initialised projections.
\item \textbf{Recipe} switches: Huginn-style random loop counts~\citep{geiping2025scaling}, truncated
  backpropagation through loops, bf16 parameters, Muon~\citep{jordan2024muon}, and the usual learning-rate
  and schedule settings.
\end{itemize}
Every switch is covered by unit tests for causality and gradient flow; the tests also verify that the
zero-initialised switches leave a warm-started model's outputs unchanged at step 0.

\subsection{Training protocol}
AdamW ($\beta=(0.9,0.95)$, weight decay 0.1, gradient clipping 1.0), peak learning rate $10^{-3}$, 100
warm-up steps and cosine decay to 5\% over the run's wall-clock budget, batch $32\times512$ bytes, fp32,
in MLX~\citep{mlx2023} on a 32\,GB Apple M4 at about 14k tokens/s. Budgets are wall clock with
compilation excluded, so a slower architecture is charged for its cost in fewer steps.

\paragraph{Two stages.} A shared \emph{base} looped model is trained once for six hours. Every experiment
loads the base, adds its modification and trains for exactly five more minutes (about 250 steps), then runs the full evaluation. The \emph{baseline} is the base plus five minutes of
unchanged training, run with three seeds; the seed spread sets the keep/discard thresholds
(\S\ref{sec:swarm}). Section~\ref{sec:calibration} explains why the protocol needs a base model.

\paragraph{Power.} Because budgets are wall clock, a laptop on battery or in Low Power Mode trains fewer
steps in the same time. Every GPU step waits for AC power with Low Power Mode off, and every ledger row
records the power state.

\subsection{Controls at matched training time}\label{sec:control-def}
To ask whether looping buys reasoning per byte, we train two conventional (non-looped) transformers
with the same code, data, optimizer and the same six-hour learning-rate schedule, and stop them after
two hours: \textbf{Plain-4L} has the same blocks and parameter count as the looped model, applied once
(effective depth 4); \textbf{Plain-10L} has the looped model's effective depth as ten distinct blocks
(about 2.4$\times$ the parameters). We compare them with the looped base's own two-hour snapshot from
the same schedule. Because time rather than steps is matched, the looped model, which does more work per
step, gets fewer optimizer steps than Plain-4L.

\subsection{Agent swarm}\label{sec:swarm-def}
Each agent is a headless Claude Code session with its own git worktree and one \emph{lane}, a family of
hypotheses described in \texttt{program.md}. An agent writes a one-line hypothesis, edits
\texttt{train.py}, and calls a runner that waits for a machine-wide GPU lock, trains, evaluates, appends a
row to a shared ledger and returns a verdict. Agents think and edit in parallel; training never
overlaps, so every run gets the same five minutes of exclusive GPU. A run is \textbf{kept} (committed to
the lane branch) if, against the lane's last kept run (or the baseline mean before the lane has one):
bpb improves by at least $\tau_{\text{bpb}}$ while OOD accuracy drops by at most $\sigma_{\text{acc}}$;
or ID or OOD accuracy improves by at least $\tau_{\text{acc}}$ while bpb worsens by at most
$\sigma_{\text{bpb}}$; or quality is equal within noise at lower cost (mean loops, KV size or
parameters). With three baseline seeds, $\tau=\max(\tau_0, 2\,\mathrm{sd})$ and
$\sigma=\max(\sigma_0,\mathrm{sd})$ with defaults $\tau_0=(0.004\text{ bpb}, 3\text{ points})$ and
$\sigma_0=(0.004, 2)$. Otherwise the agent reverts the file. After the swarm, each lane's best kept
commit is re-run with two new seeds; a lane's result counts as \textbf{confirmed} only if the three-seed
mean beats the baseline mean by more than twice the larger seed standard deviation on some metric
without losing more than that on another.
